\section{A constructive Davenport--Schmidt theorem with an explicit modulus}
\label{sec:ds}

The theorem linked in the question concerns approximation by algebraic
numbers of degree at most two, with exponent three.  It is not the later
Davenport--Schmidt theorem about simultaneous approximation to a number and
its square.  The original 1967 result gives the constant
$C>160/9$ in
\[
  |\xi-\alpha|< C\max(1,\xi^2)H(\alpha)^{-3}.
\]
Here we retain the exponent and the dependence on $\xi$, but use the
larger absolute constant $32768$.  In exchange, every search below has an
explicit primitive recursive bound.  The proof is a finite reconstruction
of the minimal-point argument of Davenport and Schmidt
\cite{DavenportSchmidt1967}, with an elementary pigeonhole estimate in
place of its sharper geometry-of-numbers estimate.

\subsection{The positive hypothesis and the exact claim}

For a nonzero integer polynomial $f=aX^2+bX+c$ of degree at most two put
$\|f\|=\max(|a|,|b|,|c|)$.  A quadratic separation modulus for $\xi$ is a
nondecreasing function $\rho:\mathbb N_{\geq1}\to\mathbb N$ such that
\begin{equation}\label{eq:ds-separation}
  0<\|f\|\leq H\quad\Longrightarrow\quad
  |f(\xi)|\geq 2^{-\rho(H)}.
\end{equation}
This is positive information.  It implies that $\xi$ is neither rational
nor quadratic.  We do not infer the existence of a primitive recursive
$\rho$ from the negative assertion that no quadratic polynomial vanishes.

\begin{theorem}[Constructive Davenport--Schmidt, with supplied separation]
\label{thm:ds-constructive}
Fix a primitive recursive rational name for a real number $\xi$, a
primitive recursive modulus $\rho$ satisfying
\eqref{eq:ds-separation}, and an integer $M\geq1$ with $|\xi|\leq M$.
For every integer $h\geq1$, a bounded primitive recursive procedure
produces a rational number or real quadratic irrational $\alpha$ such that
\[
 H(\alpha)>h,\qquad
 |\xi-\alpha|<32768\max(1,\xi^2)H(\alpha)^{-3}.
\]
The output consists of its primitive minimal polynomial and a root
selector.  All bounds used by the procedure are displayed below.

In first-order iPRA this is a schema for the fixed primitive recursive
name and modulus, conditional on their stated arithmetical correctness
properties.  It does not quantify over arbitrary functions as objects.
\end{theorem}

An algebraic real of fixed degree at least three admits such a modulus by
an elementary resultant estimate.  Thus the hypothesis has substantial
unconditional instances; it is not restricted to a hypothetical oracle.
A construction of the modulus is given at the end of this section.

We first prove a finite arithmetic core at a rational working point
$0<\eta<1$.  Temporarily, $\rho$ denotes the separation exponent used
by this core.  Separation is required only for coefficient heights at most
$W=2B^2$, where $B$ is the finite bound constructed below.  This is
consistent with $\eta$ being rational: its own linear relation may have
height much greater than $W$.  In the actual algorithm we will compute
$B$ with the exponent $\rho+1$ and obtain the working point by one
sufficiently accurate rational sample of the given real number.  Write
\[
 K=16,\qquad D=4096,\qquad C=8192=2D,\qquad
 \varepsilon(H)=2^{-\rho(H)}.
\]
The rational core has constant $C$.  Its rational sample will give the
normalized real conclusion with constant $16384$.

\subsection{A finite pigeonhole estimate}

For every integer $3\leq U\leq B$ there is a nonzero integer triple
$v=(a,b,c)$ with
\begin{equation}\label{eq:ds-dirichlet}
 \|v\|<U,\qquad 0<|a\eta^2+b\eta+c|<\frac{16}{U^2}.
\end{equation}
To prove this, put $N=\lfloor(U-1)/2\rfloor$.  The $(N+1)^2$ numbers
\[
 a\eta^2+b\eta\qquad(0\leq a,b\leq N)
\]
have fractional parts in $[0,1)$.  Partition that interval into $N^2$
half-open intervals of length $N^{-2}$.  Two distinct pairs give fractional
parts in the same interval.  Subtract them, and subtract the corresponding
integer parts.  The resulting $a,b$ have absolute values at most $N$;
the constant coefficient has absolute value at most $2N$ because the
original numbers lie in $[0,2N]$.  The triple is nonzero, since the pairs
are distinct.  Its value is nonzero by the separation hypothesis and has
absolute value less than $N^{-2}$.  Finally $2N<U$ and $N\geq U/4$.

This pigeonhole argument is constructive.  At the rational working point,
integer parts and occupied subintervals are computed by exact integer
division.  Boundary equality is decidable.  The ordinary finite
pigeonhole proof uses induction only on finite lists.  The separation
assumption is needed only for the final nonzero triple, whose height is
less than $U\leq B$.

\subsection{Finite minimal points and a bound for the next record}

For each integer $1\leq H\leq B$, consider the finite collection of nonzero
integer triples of height at most $H$.  After changing signs when needed,
their values under
\[
 L(a,b,c)=a\eta^2+b\eta+c
\]
are positive.  Choose the triple having the smallest positive value.
It is unique: equality of two values would give a nonzero quadratic
polynomial of height at most $2B\leq W$ vanishing at $\eta$.  List the distinct choices in increasing
order of the height at which they first occur, stopping at any specified
finite bound.  Denote them by $v_i$, and put
\[
 H_i=\|v_i\|,\qquad E_i=L(v_i)>0,\qquad
 P(v)=2\eta a+b,\qquad P_i=|P(v_i)|.
\]
Then $H_i$ strictly increases, $E_i$ strictly decreases, and every $v_i$
is primitive.  Indeed, a common divisor greater than one would give a
smaller positive value at a smaller height.  Consecutive $v_i$ are linearly
independent: primitive proportional integer vectors agree up to sign,
which is incompatible with their distinct positive values.

Whenever $H_{i+1}\geq3$, \eqref{eq:ds-dirichlet} applied with
$U=H_{i+1}$ gives
\begin{equation}\label{eq:ds-record-estimate}
 E_i<\frac{K}{H_{i+1}^2}.
\end{equation}
The reason is that $E_i$ is the least positive absolute value at heights
strictly below $H_{i+1}$.

Define the primitive recursive function
\begin{equation}\label{eq:ds-jump}
 J(R)=\max\{R+2,\,8\cdot2^{\rho(R)}\}\qquad(R\geq1).
\end{equation}
Whenever $J(R)\leq B$, there is a record of height strictly greater
than $R$ and strictly less than $J(R)$.  In fact, \eqref{eq:ds-dirichlet} at $J(R)$ supplies a value
less than
\[
 \frac{16}{J(R)^2}\leq\frac14\,2^{-2\rho(R)}
 <2^{-\rho(R)},
\]
whereas every nonzero triple of height at most $R$ has value at least
$2^{-\rho(R)}$.  This proves a bound for obtaining a later record without
any appeal to an infinite sequence or an unbounded search.

\subsection{The integral-plane lemma}

Two consecutive minimal points $v_i,v_{i+1}$ form an integral basis for
the integer points in their rational plane.  To see this, if an integer
point in that plane were outside their integral span, reduce its two
rational coordinates modulo integers into $[-1/2,1/2]$.  This produces
a nonzero integer point
\[
 w=r v_i+s v_{i+1},\qquad |r|,|s|\leq\tfrac12.
\]
It satisfies
\[
 \|w\|\leq\tfrac12(H_i+H_{i+1})<H_{i+1},\qquad
 |L(w)|\leq\tfrac12(E_i+E_{i+1})<E_i,
\]
contradicting minimality.  The coordinates can be computed by a nonzero
two-by-two minor, so this argument involves only rational arithmetic.

If three consecutive points are linearly dependent, the two adjacent
pairs are integral bases of the same plane.  Consequently
\begin{equation}\label{eq:ds-plane-relation}
 v_{i-1}=u v_i+s v_{i+1},\qquad u\in\mathbb Z,\quad s\in\{1,-1\}.
\end{equation}
For completeness, expansion in one integral basis gives the coefficient
$s$ as an integer.  Expansion in the other basis shows that its inverse
is an integer.  Hence $s=\pm1$.
For any two linear forms $A,B$, relation
\eqref{eq:ds-plane-relation} preserves, up to sign, the quantity
\[
 A(v_{i-1})B(v_i)-A(v_i)B(v_{i-1}).
\]
In particular, its absolute value is constant along a finite string of
dependent consecutive triples.

\subsection{A finite bound for finding an independent triple}

Suppose a pair of consecutive records $v_r,v_{r+1}$ has heights at most
$R$, and the upper bound $J(G(R))$ below is at most $B$.  Some coordinate $j$ makes
\[
 w=(v_r)_j v_{r+1}-(v_{r+1})_j v_r
\]
nonzero, because the pair is independent.  We have $\|w\|\leq2R^2$, and
therefore
\begin{equation}\label{eq:ds-wedge-lower}
 |L(w)|\geq2^{-\rho(2R^2)}.
\end{equation}
Define
\begin{equation}\label{eq:ds-gap}
 G(R)=\max\{R+1,\,32\cdot2^{\rho(2R^2)}+1\}.
\end{equation}
Take a later record $v_t$ with $H_t>G(R)$, available below $J(G(R))$.
If every triple with center $r+1,r+2,\ldots,t-1$ were dependent, the
integral-plane lemma would make the absolute value in
\eqref{eq:ds-wedge-lower} equal to
\[
 |(v_{t-1})_jE_t-(v_t)_jE_{t-1}|
 \leq H_{t-1}E_t+H_tE_{t-1}
 <\frac{2K}{H_t}
 <2^{-\rho(2R^2)}.
\]
The penultimate inequality uses \eqref{eq:ds-record-estimate} and
$E_t<E_{t-1}$.  This contradiction proves that an independent triple
occurs in the displayed finite range.  No assertion about infinitely many
independent triples has been used.

\subsection{The determinant calculation}

We now prove the key finite obstruction.  Suppose, over the finite range
under consideration, that every relevant record satisfies
\begin{equation}\label{eq:ds-small-derivative}
 P_i\leq\frac{H_i^3E_i}{D}.
\end{equation}
Let $n$ be the center of an independent triple.  The determinant with
rows
\[
 ((v_i)_1,P(v_i),L(v_i)),\qquad i=n-1,n,n+1,
\]
equals the determinant with rows $v_i$: the triangular transformation
from $(a,b,c)$ to $(a,2\eta a+b,\eta^2a+\eta b+c)$ has determinant one.
It is thus a nonzero integer and has absolute value at least one.

Expand along the middle column.  The term involving $P_{n-1}$ has absolute
value at most
\[
 P_{n-1}(H_nE_{n+1}+H_{n+1}E_n)
 <\frac{2}{D}H_{n+1}H_{n-1}^3E_{n-1}E_n
 <\frac{2K^2}{D}=\frac18.
\]
Here the last inequality follows from
$E_{n-1}<K/H_n^2$, $E_n<K/H_{n+1}^2$, and
$H_{n-1}<H_n<H_{n+1}$.  The term involving $P_n$ is similarly less than
\[
 \frac{2}{D}H_{n+1}H_n^3E_{n-1}E_n
 <\frac{2K^2}{D}=\frac18.
\]
The remaining cofactor has absolute value at most $2H_nE_{n-1}$.
Therefore
\begin{equation}\label{eq:ds-derivative-lower}
 P_{n+1}H_nE_{n-1}>\frac38,
 \qquad P_{n+1}>\frac{3H_n}{128}.
\end{equation}
The second inequality follows again from $E_{n-1}<16/H_n^2$.

Let $m<n$ be consecutive centers of independent triples.  The intervening
triples are dependent, so the integral-plane lemma gives
\begin{align*}
 \Omega&=|P(v_m)E_{m+1}-P(v_{m+1})E_m|\\
  &=|P(v_{n-1})E_n-P(v_n)E_{n-1}|.
\end{align*}
The latter expression is bounded above by
\[
 \frac1D(H_{n-1}^3+H_n^3)E_{n-1}E_n
 <\frac{2K}{D}H_nE_n=\frac{2}{256}H_nE_n.
\]
The former, by \eqref{eq:ds-derivative-lower}, is bounded below by
\[
 \frac{3}{128}H_mE_m-\frac1D H_m^3E_mE_{m+1}
 >\left(\frac{3}{128}-\frac{K}{D}\right)H_mE_m
 =\frac{5}{256}H_mE_m.
\]
We used $E_{m+1}<E_m<K/H_m^2$, so this estimate does not need a further
successor beyond $v_{m+1}$.  Combining the two bounds gives
\begin{equation}\label{eq:ds-growth}
 H_nE_n>\frac52 H_mE_m.
\end{equation}
This is a finite growth invariant, stronger than the monotonicity needed
below.

\subsection{An explicit height bound for a large derivative}

Let $T\geq16385$.  Define, in the following order,
\begin{equation}\label{eq:ds-final-bound}
\begin{aligned}
 A&=J(T),& A'&=J(A),& U&=J(G(A')),\\
 R&=\max\{U+1,16\cdot2^{\rho(U)}+1\},\\
 V&=J(R),& V'&=J(V),& B&=J(G(V')).
\end{aligned}
\end{equation}
All of these are primitive recursive functions of $T$ for the fixed
primitive recursive function $\rho$.

\begin{lemma}\label{lem:ds-large-derivative}
Among the primitive nonzero integer triples with $T<\|v\|\leq B$ there is
one satisfying
\begin{equation}\label{eq:ds-large-derivative}
 \|v\|^3|L(v)|<D|P(v)|.
\end{equation}
\end{lemma}

\begin{proof}
Assume the finite, decidable search for such a triple returns no match.
Then \eqref{eq:ds-small-derivative} holds for every record of height
strictly between $T$ and $B$.

The first record above $T$ has height below $A$; its successor has height
below $A'$.  The finite independence lemma supplies an independent triple
whose center $m$ is after that pair's first point and has $H_m<U$.
In particular, all three points involved have height greater than $T$.
The separation hypothesis and monotonicity of $\rho$ give
\[
 H_mE_m\geq E_m\geq2^{-\rho(U)}.
\]
Likewise the first record above $R$ and its successor have heights below
$V$ and $V'$.  Applying the finite independence lemma again supplies an
independent center $n$ with
\[
 R<H_n<B,
\]
whose successor is also within the searched range.  Apply
\eqref{eq:ds-growth} successively to the finite list of independent
centers from $m$ to $n$.  It follows that
$H_nE_n>2^{-\rho(U)}$.  But
\[
 H_nE_n<\frac{16}{H_n}
 <2^{-\rho(U)}
\]
by \eqref{eq:ds-record-estimate} and the definition of $R$.  This is a
contradiction.

The search is over a finite decidable list.  Its failure leads to a
contradiction, so its other, explicitly recorded outcome supplies a
witness.  This last step is finite decidable case analysis, not an
application of double-negation elimination for an unbounded existential.
\end{proof}

\subsection{Turning the large derivative into a nearby algebraic number}

Let $f=aX^2+bX+c$ be the primitive polynomial supplied by
Lemma~\ref{lem:ds-large-derivative}, and write $H=\|f\|$,
$L=f(\eta)$, $P=f'(\eta)$.  Then $L\ne0$, $P\ne0$, $H>4D$, and
\[
 |L|<D|P|H^{-3}.
\]
We claim that $f$ has a real root $\alpha$ with
\begin{equation}\label{eq:ds-root-distance}
 |\eta-\alpha|<2DH^{-3}.
\end{equation}

If $|P|\geq H^{-1}$, put $t=2DH^{-3}$ and choose
$s=-\operatorname{sgn}(LP)t$.  The exact quadratic Taylor identity gives
\[
 \operatorname{sgn}(L)f(\eta+s)
 \leq |L|-|P|t+H t^2
 <-D|P|H^{-3}+4D^2H^{-5}<0.
\]
The last inequality uses $H>4D$.  Thus the values of $f$ at $\eta$ and
$\eta+s$ have opposite signs, and its explicit linear or quadratic root
formula supplies a root strictly between them.  This use of a sign
change needs no general intermediate-value principle: it is just the
elementary quadratic formula, or equivalently its elementary algebraic
sign verification.  Indeed, throughout $|z|\leq t$,
\[
 |f'(\eta+z)-P|\leq2Ht=4D/H^2<1/H\leq|P|,
\]
so the derivative keeps the sign of $P$.  If $a\ne0$, the discriminant
$\Delta=b^2-4ac$ is positive, and the selected root is explicitly
\[
 \alpha=\frac{-b+\operatorname{sgn}(P)\sqrt{\Delta}}{2a}.
\]

If $|P|<H^{-1}$, the integer discriminant satisfies
\[
 |b^2-4ac|=|P^2-4aL|
 <H^{-2}+4DH^{-3}<1.
\]
Hence it is zero.  Here $a\ne0$, since otherwise $P=b=0$.
Primitivity implies
\[
 f(X)=\pm(uX+v)^2
\]
for coprime integers $u,v$ with $u\ne0$.  To see this, put the repeated
rational root in lowest terms $-v/u$.  Then
$f=a(X+v/u)^2$.  Since $c=av^2/u^2$ is an integer and
$\gcd(u^2,v^2)=1$, we have $a=ku^2$ for an integer $k$.  Thus the
coefficients are $ku^2,2kuv,kv^2$.  The three numbers
$u^2,2uv,v^2$ have greatest common divisor one, so primitivity forces
$k=\pm1$.  Since $u\eta+v\ne0$, the
large-derivative inequality reduces to
\[
 |u\eta+v|<2D|u|H^{-3}.
\]
The rational root $\alpha=-v/u$ again satisfies
\eqref{eq:ds-root-distance}.

In both cases $H(\alpha)\leq H$.  For an irreducible quadratic this is
equality.  For a rational root in lowest terms $p/q$, the rational-root
divisibility relations give $|p|\leq|c|$, $q\leq|a|$ when $ac\ne0$;
the cases of a zero constant coefficient and a linear polynomial are immediate.  Consequently
\[
 |\eta-\alpha|<8192H(\alpha)^{-3}.
\]

\subsection{Forcing the output height to exceed a prescribed height}

Fix $h\geq1$.  If $\beta$ has degree at most two and $H(\beta)\leq h$,
let $g$ be its primitive minimal polynomial.  For
$|\eta-\beta|<1$ we have $|\beta|<2$, and
\[
 |g(\eta)|
 =|\eta-\beta|\,|a(\eta+\beta)+b|
 \leq4h|\eta-\beta|.
\]
The same formula covers a linear $g$ by setting $a=0$.  Therefore
\begin{equation}\label{eq:ds-small-height-separation}
 |\eta-\beta|\geq\frac{2^{-\rho(h)}}{4h}.
\end{equation}
Set
\begin{equation}\label{eq:ds-threshold}
 T=\max\{16385,\,32768h\cdot2^{\rho(h)}+1\}.
\end{equation}
For every $H>T$, the root distance in
\eqref{eq:ds-root-distance} is less than both $1$ and the right hand side
of \eqref{eq:ds-small-height-separation}.  The produced root must therefore
have height greater than $h$.  Equations
\eqref{eq:ds-jump}, \eqref{eq:ds-gap}, \eqref{eq:ds-final-bound}, and
\eqref{eq:ds-threshold} are the promised explicit height-search bound.

\subsection{Freezing one rational sample: the strict iPRA version}

We now distinguish the given real $x\in(0,1)$ from the rational working
point $\eta$.  Let $\rho$ be the supplied separation exponent for $x$.
Put
\[
 \sigma(H)=\rho(H)+1.
\]
Compute $T,J,G,B$ from
\eqref{eq:ds-jump}, \eqref{eq:ds-gap}, \eqref{eq:ds-final-bound}, and
\eqref{eq:ds-threshold}, using $\sigma$ everywhere in those formulas.
Set $W=2B^2$.  Take one rational sample $\eta$ of the given name such that
\begin{equation}\label{eq:ds-freeze}
 |\eta-x|\leq
 \min\left\{\frac{2^{-\rho(W)}}{8W},\,B^{-3}\right\}.
\end{equation}
This is a primitive recursive precision request.  For example, a name
whose $s$th rational has error at most $2^{-s}$ can be sampled at
\[
 s=\rho(W)+8W+3B+10.
\]
The deliberately excessive additions avoid any need to compute logarithms.

The inequalities $x\geq2^{-\rho(1)}$ and
$1-x\geq2^{-\rho(1)}$ follow by applying separation to $X$ and $X-1$.
Consequently \eqref{eq:ds-freeze} ensures $0<\eta<1$.  For any polynomial
$f$ of height at most $H\leq W$, the elementary estimate
\[
 |f(\eta)-f(x)|\leq4H|\eta-x|
 \leq\tfrac12 2^{-\rho(W)}
 \leq\tfrac12 2^{-\rho(H)}
\]
shows
\begin{equation}\label{eq:ds-local-separation}
 |f(\eta)|\geq2^{-\sigma(H)}
 \quad(0<\|f\|\leq H\leq W).
\end{equation}
Thus the entire finite core can be run at this one rational $\eta$.
All of its separation uses lie inside this local range: differences of
searched triples have height at most $2B$; the integer vectors in the
independence lemma have height at most $2R^2$ with $R<B$; the next-record
and small-output-height arguments use heights below $B$.

The core gives a root $\alpha$ of a primitive polynomial of height
$H\leq B$, with $H(\alpha)>h$ and
$|\eta-\alpha|<8192H^{-3}$.  Therefore
\[
 |x-\alpha|
 \leq|x-\eta|+|\eta-\alpha|
 <B^{-3}+8192H^{-3}
 \leq8193H(\alpha)^{-3}
 <16384H(\alpha)^{-3}.
\]
The output-height exclusion was already proved for the rational working
point, so no separate limiting argument is needed.

This freezing step is also the foundation audit.  Every induction in
the core concerns finite integer or rational codes at the single rational
input $\eta$.  Values such as $P_iE_j-P_jE_i$ and $H_iE_i$ are rational
numbers, so their signs and inequalities are decidable by cross
multiplication; no hidden separation modulus for quartic expressions is
required.  Local condition \eqref{eq:ds-local-separation} is one bounded
primitive recursive predicate on the finite input.  Quantifier-free
induction is sufficient for the finite list invariants and bounded
searches.  Only the final error estimate refers back to the given name.

\subsection{Removing the normalization}

The separation modulus decides the signs of $\xi$, $\xi-1$, and
$\xi+1$; none is zero.  Reflection preserves algebraic heights.  Thus
if $0<|\xi|<1$, apply the normalized theorem to $x=|\xi|$ and then
reflect the output when necessary.

If $|\xi|>1$, put $x=1/|\xi|\in(0,1)$.  Reversal of the three
coefficients and multiplication by $\xi^2$ show that
\[
 |ax^2+bx+c|
 \geq M^{-2}2^{-\rho(H)}
 \geq2^{-(\rho(H)+2M)}.
\]
Thus $\rho_x(H)=\rho(H)+2M$ is an admissible primitive recursive
separation exponent for $x$.  A primitive recursive rational name
for $x$ is obtained by division, using $|\xi|>1$ as the denominator
bound.  Run the normalized construction with requested minimum height
\[
 h'=\max\{h,32768M\}.
\]
Its output $\beta$ satisfies
$|x-\beta|<16384H(\beta)^{-3}<1/(2M)$, so
$|\beta|>1/(2|\xi|)>0$.  Let $\alpha$ be the signed reciprocal of
$\beta$.  Reversal of a primitive minimal polynomial preserves its
height, and hence $H(\alpha)=H(\beta)>h$.  Finally
\[
 |\xi-\alpha|
 =\frac{|x-\beta|}{x|\beta|}
 <2\xi^2|x-\beta|
 <32768\xi^2H(\alpha)^{-3}.
\]
This completes the constructive proof of
Theorem~\ref{thm:ds-constructive}.

\subsection{Why the construction is primitive recursive}

Here are details that prevent the phrase ``effective'' from concealing
an unbounded search.  For the fixed primitive recursive $\rho$, all of
$T,J,G,B,W,s$ are explicit finite compositions of primitive recursive
functions.  After sampling once at precision $s$, enumerate the triples
in $[-B,B]^3$, test their content by Euclid's algorithm, and retain those
in the specified height range.  At this rational input every value,
minimum, comparison, record, and large-derivative test is exact rational
arithmetic.  No further refinement loop or sign oracle is used.

Quadratic root output requires only an integer square-root test and a
choice between two explicit roots.  The square root of a positive integer
has a primitive recursive rational name obtained by a prescribed number
of rational bisections.  Reduction of a rational root to lowest terms is
primitive recursive.  The root's primitive polynomial and sign selector
are finite data; approximation of its value to any requested precision
is a separate bounded primitive recursive calculation.

The iPRA proof of the rational core uses the following finite invariants:
the smallest positive value seen so far; the height and value ordering
of records; the integral-plane relation and its preserved two-by-two
minors; and the growth inequality \eqref{eq:ds-growth} on a finite list.
After multiplying by positive denominators, each is a predicate on
integer codes.  Bounded quantifiers are compiled into Boolean primitive
recursive tests before quantifier-free induction is applied.  The local
separation predicate is an explicit premise of this finite rational
theorem, and \eqref{eq:ds-freeze}--\eqref{eq:ds-local-separation} explain
how the fixed real name and modulus supply it.  The mere existence of a
terminating program would not, by itself, establish this translation;
the rational freezing is what makes its induction formulas explicit.

\subsection{An explicit separation modulus for algebraic inputs}

Suppose $\xi$ is a specified real root of an irreducible integer
polynomial $F$ of degree $d\geq3$, with coefficient height $A\geq1$ and
$|\xi|\leq R$, where $R\geq1$ is an integer.  For any nonzero integer
polynomial $g$ of degree $e\leq2$ and height at most $H$, its resultant
with $F$ is a nonzero integer.  The Sylvester-matrix adjugate supplies
integer polynomials $U,V$ with
\[
 UF+Vg=\operatorname{Res}(F,g),\qquad \deg V<d.
\]
Writing $M_H=\max(A,H,1)$, the Leibniz determinant formula bounds every
coefficient of $V$ by
\[
 (d+1)!M_H^{d+1}.
\]
Indeed the Sylvester matrix has size at most $d+2$, and its cofactors
are determinants of size at most $d+1$ with entries bounded by $M_H$.
For constant $g$, the direct lower bound $|g(\xi)|\geq1$ is even simpler.
Evaluation at $\xi$ yields
\[
 1\leq|\operatorname{Res}(F,g)|
 \leq|g(\xi)|\,d(d+1)!M_H^{d+1}(R+1)^{d-1}.
\]
Define
\[
 Q(H)=d(d+1)!\max(A,H,1)^{d+1}(R+1)^{d-1}.
\]
Then $|g(\xi)|\geq Q(H)^{-1}\geq2^{-Q(H)}$.
Thus $\rho(H)=Q(H)$ is an explicit nondecreasing primitive recursive
separation exponent.  The estimate is deliberately coarse.

For a fixed algebraic root one may use its usual primitive recursive
isolating-interval name.  To avoid concealing a termination assertion,
the input can instead include a rational interval whose endpoint
values of $F$ have opposite signs, together with a positive lower bound
on $|F'|$ throughout the interval; sufficiently small such intervals
exist for a simple algebraic root, and once supplied, bisection has an
explicit precision bound.  These finite certificates are checkable by
rational polynomial estimates.  Uniform construction of all algebraic
root data can be treated separately by a bounded root-isolation algorithm.

\subsection{Lean implementation plan and the remaining formal work}

The mathematical proof above is complete for the stated positive
hypotheses; no assertion is made that the entire iPRA derivation has
already been machine checked.  A suitable Lean development has these
modules, in dependency order:
\begin{enumerate}
 \item Integer triples, content, height, $L$, $P$, polynomial reversal,
       and finite list enumeration.  State the separation hypothesis for
       fixed encoded primitive recursive functions, and the bounded
       local separation predicate for a rational working point.
 \item One rational freezing operation with the explicit precision $s$,
       proving the bounded local separation predicate.  All later
       decisions use exact rational arithmetic.
 \item The finite pigeonhole lemma and the record algorithm, proving
       \eqref{eq:ds-record-estimate} and the next-record bound $J$.
 \item The integral-plane lemma using two-by-two minors and rational
       coordinate reduction; then the finite independence bound $G$.
 \item The determinant identity and the rational inequalities leading
       to \eqref{eq:ds-growth}.  These are fixed-degree polynomial
       identities and ordered rational arithmetic after denominators
       are cleared.
 \item The finite search bound $B$, proved by the two bounded
       applications of the independence lemma and finite induction over
       the list of independent centers.
 \item Linear/quadratic root certificates, the discriminant-zero
       factorization, minimal-polynomial height comparison, and the
       exclusion of heights at most $h$.
 \item Reflection and reciprocal transformations, and optionally the
       Sylvester-adjugate separation bound for algebraic inputs.
\end{enumerate}
Proving the corresponding assertions over Lean's real numbers is useful
for checking the mathematics, but it is not the final foundation audit.
The iPRA layer must prove the bounded algorithms' rational certificate
properties in the repository's syntax and exhibit the relevant
quantifier-free induction instances.  This is a clear, finite proof
engineering task; an appeal to classical real-number minimum selection
would lose the constructive content established here.
